% =====================================================================================
% sections/s3_methods.tex -- Methods.
% All `% src:` paths are relative to the root of the repository.
% Labels defined here and used elsewhere:
%   s3_methods:sec, s3_methods:eq:pinn, s3_methods:eq:ritz, s3_methods:eq:ritz-mc,
%   s3_methods:eq:errors, s3_methods:tab:setup, s3_methods:tab:samplers,
%   s3_methods:prop:robin, s3_methods:rem:consistent, s3_methods:sec:* (subsections)
% Details in the Supplementary Material: sections/S2_methods.tex (figure of the point
% sets, Observation s3_methods:obs:robin, the two Laplacians, implementation).
% =====================================================================================
\section[Methods]{Methods: PINN loss, Deep Ritz energy, collocation strategies, optimisers
and error measures}
\label{s3_methods:sec}

This section fixes what is common to the experiments of
Sections~\ref{s4_lowdim:sec}--\ref{s6_highdim:sec}. Further details are in Supplementary
Section~\ref{S2:sec}, and the settings of every experiment in
Table~\ref{sC_details:tab:hyper}.

% -------------------------------------------------------------------------------------
\subsection{Networks and the PINN loss}\label{s3_methods:sec:pinn}

Every approximant $\unet$ is a fully connected network with $\tanh$ activations and a
linear output layer, with the default initialisation of PyTorch~\citep{paszke2019},
un-normalised inputs and single-precision arithmetic (Table~\ref{s3_methods:tab:setup}).

% src: research_benchmark/pinnbench/problems.py (hidden, blocks(), terms(), pde() of each class)
% src: research_benchmark/results/timing_single.json (n_params, n_r)
% src: research_highdim/scripts/hd_core.py (MLP, train: n_int = n_bd = 1024); research_highdim/results/cost_vs_d.json (n_params 8577, 8641, 8769, 9089, 9729 at d = 2, 3, 5, 10, 20, i.e. 64d + 8449)
\begin{table}[tbp]
\centering
\caption{Network, interior budget $\Nr$ and the terms of the PINN
loss~\eqref{s3_methods:eq:pinn} for each problem. ``$2\times128$'' means two hidden layers
of 128 units. Each constraint block is one mean squared term; the number of points of a
block is given in brackets. The initial displacement and the initial velocity of the wave
problems are imposed at the same points.}
\label{s3_methods:tab:setup}
\small
\setlength{\tabcolsep}{4.5pt}
\begin{tabular}{@{}llrrl>{\raggedright\arraybackslash}p{0.37\textwidth}@{}}
\toprule
Problem & Hidden & Parameters & $\Nr$ & Residual $r_\theta$ & Constraint blocks \\
\midrule
\PH     & $2\times128$ & 17\,025 & 256  & $\partial_t\unet-\partial_{xx}\unet$
        & $u(0,\cdot)$ [100]; $u(\cdot,0)$, $u(\cdot,1)$ [100 each] \\
\PWone  & $3\times64$  & 8577    & 1024 & $\partial_{tt}\unet-4\,\partial_{xx}\unet$
        & $u(0,\cdot)$, $\partial_t u(0,\cdot)$ [100]; $u(\cdot,0)$, $u(\cdot,1)$ [100 each] \\
\PWtwo  & $2\times50$  & 2801    & 512  & $\partial_{tt}\unet-\lap\unet$
        & $u(0,\cdot)$, $\partial_t u(0,\cdot)$ [$21\times21$]; $u$ on the 4 side faces
          [$11\times21$ each] \\
\PLtwo  & $3\times128$ & 33\,537 & 256  & $\lap\unet$
        & $u$ at $x=0$, $x=\pi$ [100 each]; $\partial_y u$ at $y=0$, $y=\pi$ [100 each] \\
\PLthree, \PLthreeS & $2\times100$ & 10\,601 & 512 & $\lap\unet$
        & $u$ on the 6 faces [$12\times12$ each] \\
\Pone, \Ptwo & $3\times64$ & $64d+8449$ & 1024 & $\lap\unet+f$
        & $u$ on $\dOm$ [1024], weight $\lam$ \\
\bottomrule
\end{tabular}
\end{table}

A PINN~\citep{raissi2019} is trained by minimising the mean squared residual of the
differential equation at $\Nr$ interior points $\x_i$ plus one mean squared term for each
of the $K$ blocks of initial or boundary data,
\begin{equation}\label{s3_methods:eq:pinn}
\Lpinn(\theta)=\frac{1}{\Nr}\sum_{i=1}^{\Nr} r_\theta(\x_i)^2
+\lam\sum_{k=1}^{K}\frac{1}{N_k}\sum_{j=1}^{N_k}
\Bigl(B_k\unet\bigl(\boldsymbol{y}^{(k)}_j\bigr)-g_k\bigl(\boldsymbol{y}^{(k)}_j\bigr)\Bigr)^2 .
\end{equation}
Here a block is one edge, one face or one initial condition, with its points
$\boldsymbol{y}^{(k)}_j$, its datum $g_k$ and its operator $B_k$ (the identity,
$\partial_t$ for an initial velocity, $\partial_y$ for a Neumann edge); for the evolution
problems the points are space--time points. Because every block is a mean, blocks with
different numbers of points weigh the same. In Sections~\ref{s4_lowdim:sec}
and~\ref{s5_hard:sec} we take $\lam=1$ and tune nothing; in Section~\ref{s6_highdim:sec}
there is one boundary block and $\lam$ is chosen as described in
Section~\ref{s3_methods:sec:highdim}.
% src: research_benchmark/pinnbench/problems.py (gradient(), second())
All derivatives in $r_\theta$ and $B_k\unet$ are computed by reverse-mode automatic
differentiation of the network output with respect to the \emph{full} input tensor, and
the required component is selected by indexing the result. Differentiating with respect
to a slice of the input is not equivalent: the slice is a new tensor on which the output
does not depend (Section~\ref{s7_pitfalls:sec:autograd}).

% -------------------------------------------------------------------------------------
\subsection{Collocation strategies}\label{s3_methods:sec:colloc}

The interior points $\x_i$ are placed by one of the six strategies of
Table~\ref{s3_methods:tab:samplers}, always with the same number $\Nr$ of points. The
budgets $\Nr=256=16^2$ (\PH, \PLtwo), $1024=32^2$ (\PWone) and $512=8^3$ (\PWtwo,
\PLthree) are perfect powers, as a tensor grid requires, and powers of two, as suits a
Sobol' set. They are deliberately small, so that the placement of the points can matter;
Section~\ref{s4_lowdim:sec} also varies $\Nr$ for \PH. The constraint points are identical
for every strategy: 100 equispaced points on each edge or initial line of the two-dimensional
domains, a $21\times21$ initial grid and $11\times21$ points on each side face for \PWtwo,
and an open $12\times12$ cell-centred grid on each face for \PLthree, which keeps away
from the two edges where its data are discontinuous.
% src: research_benchmark/pinnbench/problems.py (blocks())

% src: research_benchmark/pinnbench/samplers.py (GridSampler, RandomSampler, ResampleSampler, SobolSampler, RADSampler: period=500, pool_factor=8, k=1, c=1)
% src: research_benchmark/pinnbench/core.py (run_batched: resample and RAD redraw at the branch point)
% src: verify_research_benchmark/vpinn.py (points(kind="nodegrid")); notes/VERIFICATION.md (sec. 2.1, grid-n control)
\begin{table}[tbp]
\centering
\caption{Collocation strategies for the $\Nr$ interior points. ``Coordinate'' includes
time for the evolution problems; all sets are mapped affinely to the domain. A seed fixes
the sets of the four random strategies.}
\label{s3_methods:tab:samplers}
\small
\renewcommand{\arraystretch}{1.15}
\begin{tabular}{@{}l>{\raggedright\arraybackslash}p{0.56\textwidth}
  >{\raggedright\arraybackslash}p{0.25\textwidth}@{}}
\toprule
Strategy & Construction & During training \\
\midrule
\strat{grid-c} & cell-centred tensor grid: with $D$ coordinates and $q=\Nr^{1/D}$, the
  points $(j+\tfrac12)/q$, $j=0,\dots,q-1$, in each coordinate; no point on the boundary
  or on the initial line & fixed \\
\strat{grid-n} & node-centred closed tensor grid: the points $j/(q-1)$, $j=0,\dots,q-1$,
  in each coordinate; includes the boundary and the initial line. Used only in the controls
  of Section~\ref{s4_lowdim:sec:strategy} & fixed \\
\strat{random} & independent uniform points, drawn once & fixed \\
\strat{resample} & independent uniform points & redrawn before every \Adam\ iteration;
  one further draw, then fixed, for \LBFGS \\
\strat{Sobol} & first $\Nr$ points of a scrambled Sobol'
  sequence~\citep{sobol1967,owen1998}, generated with SciPy~\citep{virtanen2020} & fixed \\
\strat{RAD} & residual-based adaptive distribution of \citet{wu2023}: starting from the
  \strat{Sobol} set, $\Nr$ points are drawn without replacement from a uniform pool of
  $8\Nr$ points with probability proportional to
  $\abs{r_\theta}/\operatorname{mean}\abs{r_\theta}+1$
  & redrawn every 500 \Adam\ iterations and once before \LBFGS \\
\bottomrule
\end{tabular}
\end{table}

% src: data/s3_methods_points.json (summary: min_t, n_points_with_t_below_1_over_32)
For \PH\ the cell-centred grid has its first column of points at $t=1/32$ and so leaves the
strip $0<t<1/32$ next to the initial line without a residual point; 14 of the 256
\strat{random} points and 8 of the 256 \strat{Sobol} points of seed 0 lie in that strip. The
node-centred grid has a column on the initial line itself but none inside the strip: its
next column is at $t=1/15$ (Figure~\ref{s3_methods:fig:points}). This matters in
Section~\ref{s4_lowdim:sec:strategy}.

% -------------------------------------------------------------------------------------
\subsection{Optimiser arms}\label{s3_methods:sec:optim}

Two arms of 3000 iterations each are compared, an iteration being one update of the
parameters. The arm \Adam\ runs Adam~\citep{kingma2015} with the constant learning rate
$10^{-3}$ for 3000 iterations. The arm \AdamLBFGS\ shares the first 1500 of these
iterations---it branches from a stored copy of the same run---and continues with 1500
iterations of L-BFGS~\citep{liu1989} with memory 50 on a fixed point set, a common
schedule~\citep{lu2021}.
% src: research_benchmark/pinnbench/core.py (ADAM_LR, LBFGS_MEMORY, run_single, run_batched)
% src: research_benchmark/results/timing_single.json (lbfgs_fevals_per_iter: 1.06, 1.06, 1.065, 1.045, 1.055)
For a single network, with the L-BFGS of PyTorch and a strong-Wolfe line search, L-BFGS needed
between 1.045 and 1.065 loss-and-gradient evaluations per iteration on the five problems, so
an equal number of iterations is close to an equal cost. The multi-seed runs use our own
L-BFGS for a stack of networks (Supplementary Section~\ref{S2:sec:impl}); the comparison of the
arms is therefore one at equal iterations.

% -------------------------------------------------------------------------------------
\subsection{The Deep Ritz method and its boundary penalty}\label{s3_methods:sec:ritz}

For the Dirichlet problem $-\lap u=f$ in $\Om=(0,1)^d$, $u=g$ on $\dOm$, the Deep Ritz
method~\citep{eyu2018} minimises over the network parameters the energy whose
Euler--Lagrange equation is the differential equation, with the boundary condition
imposed by a penalty of weight $\bet$:
\begin{equation}\label{s3_methods:eq:ritz}
\Eritz(v)=\int_{\Om}\Bigl(\tfrac12\abs{\grad v}^2-fv\Bigr)\dd\x
+\bet\int_{\dOm}(v-g)^2\,\dd s .
\end{equation}
With $\Nr$ points $\x_i$ uniform in $\Om$ and $\Nb$ points $\boldsymbol{y}_j$ uniform on
$\dOm$, and because $\abs{\Om}=1$ and $\abs{\dOm}=2d$, the loss is the Monte-Carlo estimate
\begin{equation}\label{s3_methods:eq:ritz-mc}
\frac{1}{\Nr}\sum_{i=1}^{\Nr}\Bigl(\tfrac12\abs{\grad\unet(\x_i)}^2-f(\x_i)\,\unet(\x_i)\Bigr)
+\frac{2d\,\bet}{\Nb}\sum_{j=1}^{\Nb}\bigl(\unet(\boldsymbol{y}_j)-g(\boldsymbol{y}_j)\bigr)^2 .
\end{equation}
% src: research_highdim/scripts/hd_core.py (loss_ritz, sample_boundary)
Only first derivatives of the network enter. Unlike the PINN loss, the minimum of this
loss is not zero and is not known in advance, so its raw value is not an error indicator:
only its excess over the minimum measures an error, namely
$\Eritz(v)-\Eritz(\ubeta)=\tfrac12\norm{\grad(v-\ubeta)}_{L^2(\Om)}^2+\bet\,\norm{v-\ubeta}_{L^2(\dOm)}^2$,
where $\ubeta$ is the minimiser of $\Eritz$. The distance measured is that to $\ubeta$, not
to the solution of the Dirichlet problem.

The penalty changes the problem that is solved. This is classical~\citep{babuska1973} and
has been analysed for the Deep Ritz method by \citet{muller2022}. We claim no novelty for
the following statement; its proof on the cube is elementary and is included to keep the
article self-contained (Supplementary Section~\ref{sA_proofs:sec:robin}).

\begin{proposition}[Penalised Ritz energy: Robin problem and penalty bias]%
\label{s3_methods:prop:robin}
\StmtRobin
\end{proposition}

\noindent On the cube the constant of part~(c) is sharpened by
Corollary~\ref{s6_highdim:cor:bound}(b) in Section~\ref{s6_highdim:sec:bound}.
% src: verify_research_highdim/results/v_robin_fd.json (rel_l2_N200: 5.00e-2, 5.48e-3, 5.54e-4 at beta = 10, 100, 1000; norm of u* = 1/3)
For \Pone\ at $d=2$, a finite-difference solution of the Robin problem shows the relative bias
$\norm{\ubeta-\uex}/\norm{\uex}$ to decay like $0.50/\bet$ to $0.55/\bet$ for $\bet\ge10$,
that is, the bias $\norm{\ubeta-\uex}$ itself like $0.167/\bet$ to $0.185/\bet$
(Observation~\ref{s3_methods:obs:robin}). The bias
is the error of the limit of the training, not a lower bound on the error of a trained
network, which can end on either side of it.

\begin{remark}\label{s3_methods:rem:consistent}
\StmtConsistent
\end{remark}

% -------------------------------------------------------------------------------------
\subsection{Protocol in \texorpdfstring{$d$}{d} dimensions}\label{s3_methods:sec:highdim}

For \Pone\ and \Ptwo\ the PINN and the Deep Ritz method use the same network, optimiser
and points and differ only in the loss: \eqref{s3_methods:eq:pinn} with residual
$\lap\unet+f$ and one boundary block of weight $\lam$,
or~\eqref{s3_methods:eq:ritz-mc}. At every iteration $\Nr=1024$ interior and $\Nb=1024$
boundary points are drawn afresh, as in~\citep{sirignano2018}; a boundary point is
obtained by choosing one of the $2d$ faces uniformly. The optimiser is Adam with learning
rate $10^{-3}$, multiplied by $0.1$ after 50 and after 75 per cent of the run, for 4000
iterations unless stated otherwise. For a given seed the two methods start from the same
network and see the same points.
% src: research_highdim/scripts/hd_core.py (train, sample_interior, sample_boundary)
The weights were chosen once for each problem and method, from $\{1,10,100,1000\}$, at
$d=5$, by the error of one run with a separate seed on a validation set of 5000 points:
$\lam=1000$ for the PINN on both problems, $\bet=100$ for Deep Ritz on \Pone\ and $\bet=1$
on \Ptwo. They are then used for every $d$. Because $\lam$ weights a boundary mean and $\bet$
a boundary integral (Remark~\ref{s3_methods:rem:consistent}), keeping both fixed means that
the coefficient of one boundary sample in the Deep Ritz loss is $2d\,\bet/\lam$ times its
coefficient in the PINN loss: 0.4 at $d=2$ and 4 at $d=20$ on \Pone, and 0.004 and 0.04 on
\Ptwo. The comparison across $d$ in Section~\ref{s6_highdim:sec} is made under this
convention; a convention that fixed the weight per unit of boundary measure for both
methods was not tried.
% src: research_highdim/scripts/hd_core.py (loss_pinn, loss_ritz: boundary terms); research_highdim/results/sweep.csv
The Laplacian in the PINN residual is computed by propagating $(u,\grad u,\lap u)$ forward
through the layers~\citep{li2024forward}; nested reverse-mode differentiation gives the same
loss at a higher cost (Supplementary Section~\ref{S2:sec:lap}).

% -------------------------------------------------------------------------------------
\subsection{Error measures and trivial predictors}\label{s3_methods:sec:errors}

Errors are measured on the held-out evaluation points
$\boldsymbol{z}_1,\dots,\boldsymbol{z}_M$ of Table~\ref{s2_problems:tab:problems}, and
always for the network at the last iteration: no checkpoint is selected. With
$\norm{v}_2^2=\sum_{i=1}^{M}\abs{v(\boldsymbol{z}_i)}^2$ and
$\overline{\uex}=M^{-1}\sum_{i}\uex(\boldsymbol{z}_i)$,
\begin{equation}\label{s3_methods:eq:errors}
\begin{gathered}
\errL=\frac{\norm{\unet-\uex}_2}{\norm{\uex}_2},\qquad
\errI=\max_{1\le i\le M}\abs{\unet(\boldsymbol{z}_i)-\uex(\boldsymbol{z}_i)},\qquad
\errC=\frac{\norm{\unet-\uex}_2}{\norm{\uex-\overline{\uex}}_2},\\[0.4em]
\errH=\frac{\norm{\grad\unet-\grad\uex}_2}{\norm{\grad\uex}_2},\qquad
\errT(t)=\frac{\norm{\unet(t,\cdot)-\uex(t,\cdot)}_{2,t}}{\norm{\uex(t,\cdot)}_{2,t}},
\end{gathered}
\end{equation}
where $\norm{\cdot}_{2,t}$ sums over the evaluation points of the time slice $t$ only.
% src: research_benchmark/pinnbench/core.py (error_metrics); research_highdim/scripts/hd_core.py (evaluate)
We report $\errL$ and $\errI$ for every problem; the centred error $\errC$ and the
$H^1$-seminorm error $\errH$ in Section~\ref{s6_highdim:sec}; and the slice error
$\errT$ for \PH. A relative error has to be read against what a trivial predictor achieves:
the zero function has $\errL=1$, the best constant has $\errC=1$, and
Proposition~\ref{s2_problems:prop:floors} gives the errors of the best constant and of the
best affine function for \Pone\ and \Ptwo; for \Pone\ they decrease with $d$, so a small
$\errL$ in high dimension is not by itself evidence of an accurate solution.

% -------------------------------------------------------------------------------------
\subsection{Seeds and statistical procedure}\label{s3_methods:sec:stats}

A seed fixes the initial weights and every random point set of a run. We use 10 seeds for
\PH\ and \PLtwo, 5 for \PWone, \PWtwo\ and \PLthree, and 3 in
Section~\ref{s6_highdim:sec}; other counts are stated where they occur. A \emph{network} is
one combination of problem, strategy and seed, trained up to the iteration at which the
optimiser arms branch; a \emph{run} is one network in one arm; a \emph{pair} is the two runs
of one network. The benchmark thus has 175 networks, 175 pairs and 350 runs. The strategies
of one seed start from the same network and are compared \emph{paired by seed}. We report
means with sample standard deviations over seeds, and geometric means; the ratio of two
strategies or arms is the geometric mean of the per-seed ratios.
% src: research_benchmark/scripts/analyze.py

The test is the exact two-sided Wilcoxon signed-rank test~\citep{wilcoxon1945} on the
per-seed differences of $\log\errL$. With $n$ seeds its smallest attainable $p$-value is
$2^{1-n}$, reached when all seeds agree in direction: $2/1024\approx0.002$ for 10 seeds and
$0.0625$ for 5.
% src: data/s4_lowdim_numbers.json (optimiser.*.per_seed_wilcoxon_p: 0.001953125 with 10 seeds, 0.0625 with 5)
For the comparison of the optimiser arms the errors of the five strategies are first
averaged geometrically within each seed, because they share an initial network. The
$p$-values are not corrected for multiple comparisons. We call an effect
\emph{significant} only if it has the smallest attainable ten-seed value, $p=2/1024$, that is,
only if all ten seeds agree in direction.
% src: sections/S8_details.tex (Table sC_details:tab:paired: 16 ten-seed strategy tests, 2 ten-seed optimiser tests); data/s4_lowdim_gridcontrol.json (6 ten-seed tests printed in Table s4_lowdim:tab:gridcontrol and its text)
There are 24 signed-rank tests with ten seeds in the article (16 strategy comparisons,
2 optimiser comparisons and 6 tests of the control experiment of
Section~\ref{s4_lowdim:sec:strategy}); by the Bonferroni inequality the probability that one
or more of them is called significant when no effect exists is at most
$24\times2/1024=4.7$ per cent. Nothing measured with fewer than ten seeds can be called
significant by this rule; for those comparisons we report directions and ranges, and their
$p$-values are descriptive. With 3 seeds (Section~\ref{s6_highdim:sec}) no test is made, and
statements are restricted to cases in which the seed ranges of the two methods are disjoint.
The rule was fixed after the results were known; it is stated so that it can be judged, not
because it was pre-registered. The looser threshold $p\le0.004$ would also admit effects on
which nine of ten seeds agree and the one dissenting seed has the smallest difference (signed-rank
statistic $W=1$, $p=4/1024$); two such comparisons exist, Sobol' points on \PLtwo\ and the
repeated comparison of the cell-centred grid with random points in the control experiment of
Section~\ref{s4_lowdim:sec:strategy}, and we call them suggestive, not significant.

% -------------------------------------------------------------------------------------
\subsection{Finite-difference baselines}\label{s3_methods:sec:fd}

Each low-dimensional problem is also solved, on a sequence of grids, by a standard
finite-difference scheme of second order~\citep{leveque2007}: the Crank--Nicolson
scheme~\citep{crank1947} with $\Delta t=\Delta x$ for \PH; the explicit leapfrog scheme
with Courant number $0.5$ and a second-order starting step for \PWone\ and \PWtwo\ (for
\PWone\ the Courant number 1 is avoided because the scheme is then exact at the nodes,
which would not be a representative baseline); the five-point Laplacian with second-order
ghost-point Neumann rows and a sparse direct solver for \PLtwo; and the seven-point
Laplacian, solved with a discrete sine transform, for \PLthree. The grid solution is
evaluated on the same held-out points as the networks, by its nodal values where these
points are nodes and by cubic spline interpolation otherwise.
% src: research_benchmark/pinnbench/classical.py

% -------------------------------------------------------------------------------------
\subsection{Verification codes and pre-registration}\label{s3_methods:sec:verif}

Two bodies of code, written separately from the code of the experiments, are used to check
the results, and are referred to by name.
% src: verify_research_benchmark/vpinn.py; verify_research_highdim/scripts/v_core.py; notes/VERIFICATION.md (sec. 2.1; sec. 3.1, V10); verify_research_highdim/results/v_summary.txt (seeds per cell)
\textsf{V-bench} re-implements the low-dimensional PINNs: one network at a time on the
CPU, derivatives propagated analytically through the layers, the strong-Wolfe \LBFGS\ of
PyTorch, seeds 0 to 5. It re-ran 40 configurations of
Sections~\ref{s4_lowdim:sec} and~\ref{s5_hard:sec}. \textsf{V-dim} re-implements the
$d$-dimensional study: its own exact solutions, a nested Laplacian, its own test points,
seeds 7 to 9 (7 to 11 in one cell); it trained 73 networks for Section~\ref{s6_highdim:sec}.
Both also re-aggregated the stored raw results (Supplementary
Section~\ref{sC_details:sec:verification}).

The verification has two limits. First, the codes were written, and the stored results
re-aggregated, separately from the experiments but within the same project, with the aim of
refuting the results; they ran on the same machine, with the same library versions, as the
experiments. This guards against errors of one code and of one analysis; it is not a
replication by another group on other hardware, and errors common to both sides (a shared
misreading of a problem, the single-precision library stack) are not excluded. Second, the
following experiments were run with the code of the experiments only and are not covered:
the longer runs on \PWtwo\ and \PWone\ (Sections~\ref{s5_hard:sec:w2}
and~\ref{s5_hard:sec:w1}), the compatible-data problem \PLthreeS\
(Section~\ref{s5_hard:sec:l3}), the slice errors of \PH\ (Section~\ref{s4_lowdim:sec:timeslice}),
the control experiment for the cell-centred grid (Section~\ref{s4_lowdim:sec:strategy}), the
16\,000-iteration runs at $d=20$ (Section~\ref{s6_highdim:sec:budget}) and the demonstration
of Section~\ref{s7_pitfalls:sec:autograd}. The demonstration of Section~\ref{s7_pitfalls:sec:oneface}
was re-run for three of its five seeds with the network and derivative code of \textsf{V-bench}.
% src: data/s4_lowdim_gridcontrol.json (agreement_with_benchmark: 0.944, 0.967); data/s4_lowdim_timeslices.json (space-time geometric means 1.07e-3, 1.16e-3 against 1.19e-3, 1.15e-3); data/s5_hard_wave2d_long.json and data/s5_hard_wave1d_long.json (errors at iteration 3000 against the benchmark runs of the same seeds); data/sC_details_repro.json (bit-for-bit re-runs)
The experiments listed above have only internal cross-checks: the configurations that they share with the benchmark
agree with it, \PLthreeS\ is run with its own control, and the drivers are reproducible bit
for bit on the machine used (Supplementary Section~\ref{sC_details:sec:repro}).

The analyses of Sections~\ref{s6_highdim:sec:bound} and~\ref{s6_highdim:sec:remedies} were
pre-registered: hypotheses, networks or problems, seeds and decision rules were fixed in a
document whose hash and time were recorded in the repository (\texttt{PREREG\_FREEZE.json}) before
the first counted evaluation or run; no external registry was used. Each was
then re-checked, in two rounds, with separately written code within the same project and on the
same machine: the proofs were re-read and their formulas checked numerically, every quoted number
was recomputed from the stored files, networks were retrained with other seeds, and problems were
added on which neither remedy had been tried. Below, the two rounds are referred to together as
\emph{the re-check}; it has the limits stated above; where it added a result, the text says so, and
Supplementary Section~\ref{sC_details:sec:verification} describes it.
% src: research_highdim_bound/PREREG_T2.md; research_highdim_bound/PREREG_FREEZE.json; research_highdim_remedies/PREREG_C1.md; research_highdim_remedies/PREREG_C1_P5.md; research_highdim_remedies/PREREG_FREEZE.json; notes/VERIFICATION.md (sec. 4)
